\chapter*{Preface}
\addcontentsline{toc}{chapter}{Preface}

CNA is a C\texttt{++}23 reimplementation of the Microsoft XNA~4.0 programming model. It uses
SDL3 for platform integration and selects one graphics renderer at build time. At this edition's
pin, the public selector exposes \RendererIdentityCount{} renderer identities implemented by
\RendererFamilyCount{} source families. CNA is neither a game nor an editor-centered engine. It
is a compatibility framework for native C\texttt{++} programs written in the shape of
\texttt{Microsoft::Xna::Framework}.

The surrounding repositories provide capabilities CNA should not duplicate.
\texttt{sharp-runtime} supplies selected .NET-shaped types; \texttt{easy-gl} and
\texttt{meta-gl} provide the OpenGL/OpenGL~ES layer used by five public identities;
\texttt{free-direct} and \texttt{free-api} implement a narrow legacy compatibility route;
\texttt{xna4-spec} makes the XNA API surface searchable; and the sample and example repositories
exercise the framework from a user's perspective. Chapter~\ref{ch:ecosystem-vocabulary} maps
those relationships without treating every adjacent repository as part of CNA itself.

This book serves two roles. It is a technical monograph about CNA's architecture and engineering
trade-offs, and a reference for people building, porting, testing, or extending CNA programs.
It follows public API contracts into their pinned implementations, then separates what source
inspection establishes from what a build, execution, renderer trace, independent oracle, or
pixel comparison establishes. Bugs, partial routes, stale project documents, and unsupported
configurations remain visible because they affect engineering decisions.

The book is not the project's live API documentation. Software changes faster than an edition
can. Treat the text as a revision-bounded map: use it to locate the relevant subsystem and its
constraints, then check the current repository before relying on a detail in new code. The full
pin set and observed checkout state appear in Appendix~\ref{ch:repo-map}.

The organization follows the path of a call through the system. Parts~\ref{part:ecosystem}
and~\ref{part:architecture-framework} introduce the ecosystem, build, framework lifecycle,
math conventions, and module boundaries. Parts~\ref{part:graphics-machine}
and~\ref{part:renderers} trace the public graphics API into renderer implementations and native
APIs. Parts~\ref{part:content-pipeline} and~\ref{part:3d-content} cover asset resolution, XNB,
CNJ, glTF, models, animation, and conformance evidence. Part~\ref{part:services} covers input,
audio, media, devices, storage, networking, and related XNA services. The remaining parts cover
foundation libraries, platforms, verification, and porting practice.

Appendix~\ref{ch:verification-tiers} is the authoritative evidence key. In short, API presence,
successful compilation, execution, renderer engagement, behavioral comparison, and pixel
comparison are different observations. A green result is reported at the narrowest level the
retained evidence supports. This convention lets a reader distinguish a working portable path
from a source-visible implementation, a historical run, or a test that skipped before reaching
its intended boundary.

The production process, including its use of AI-agent assistance, is described in the front
matter that follows. The authority for technical claims remains the identified source and
evidence, not the tool that drafted a sentence.
